\documentclass[11pt,a4paper]{article}
\usepackage{shared/luxcover}
\usepackage[utf8]{inputenc}
\usepackage[T1]{fontenc}
\usepackage{amsmath,amssymb,amsfonts,amsthm}
\usepackage{graphicx}
\usepackage{hyperref}
\usepackage{booktabs}
\usepackage{algorithm}
\usepackage{algpseudocode}
\usepackage{enumitem}
\usepackage{xcolor}

\hypersetup{colorlinks=true,linkcolor=black,citecolor=blue,urlcolor=blue}

\newtheorem{definition}{Definition}[section]
\newtheorem{remark}{Remark}[section]

\title{LP-306: LX --- Decentralized Social Trading Exchange\\[4pt]
\large From the 2022 Whitepaper to the Lux DEX Production Stack}

\author{
Zach Kelling\thanks{Corresponding author: zach@lux.network}\\
\textit{Lux Industries, Inc.}\\
\texttt{research@lux.network}
}

\date{Original: August 2022\quad|\quad Revised: March 2026}

\begin{document}

\luxcoverpage

\begin{abstract}
LX was proposed in 2022 as a decentralized social trading platform
combining copy trading, people-based portfolios, atomic swaps, and
a DAO-governed liquidity protocol on the Lux Network. This paper
preserves the original design, maps each component to its current
implementation status across the Lux DEX
(\texttt{luxfi/dex}), Lux CEX (\texttt{luxfi/cex}), Lux Broker
(\texttt{luxfi/broker}), and the Lux ATS (\texttt{luxfi/ats}), and
identifies the features that shipped, evolved, or were superseded by
superior alternatives.
\end{abstract}

\tableofcontents
\newpage

%% ============================================================
\section{Introduction}

The 2022 LX whitepaper described a social trading DEX with the
following thesis: centralized exchanges are single points of failure
(citing MtGox, Bitfinex, Coincheck); decentralized alternatives
at the time (EtherDelta, 0x Protocol) suffered from latency, poor
UX, and front-running. LX proposed to solve both with a non-custodial,
zero-knowledge encrypted orderbook on a high-throughput L1.

Four years later, the Lux Network ships three exchange-layer
components that realize this vision:

\begin{center}
\begin{tabular}{llr}
\toprule
\textbf{Component} & \textbf{Repo} & \textbf{Throughput} \\
\midrule
Lux DEX (native orderbook) & \texttt{luxfi/dex} & 434M ops/sec (GPU) \\
Lux CEX (compliance gateway) & \texttt{luxfi/cex} & FIX 4.4 + REST + WS + ZAP \\
Lux Broker (16-venue SOR) & \texttt{luxfi/broker} & Smart order routing \\
\bottomrule
\end{tabular}
\end{center}

%% ============================================================
\section{Architecture: 2022 Design vs.\ 2026 Implementation}

\subsection{Matching Engine}

\textbf{2022}: C++ matching engine compiled to Solidity WASM,
running on a Lux subnet. Order book stored on-chain.

\textbf{2026}: Multi-engine architecture
(\texttt{luxfi/dex}): Go (default), C++ (high-freq), Rust
(safety-critical), MLX (Apple Silicon). CLOB matching with
AMM fallback. Clearinghouse with margin and liquidation.
Cross-chain gateway bridging 8 chains. 434M ops/sec on GPU
via the NTT-accelerated matching kernel (LP-203).

\subsection{Social Trading and Copy Swaps}

\textbf{2022}: Permission-mapped liquidity aggregator smart
contracts allowing expert traders to manage investor capital
on-chain while investors retain key custody. People-Based
Portfolios (PBPs) with transparent performance stats.

\textbf{2026}: The Lux ATS (\texttt{luxfi/ats}) implements the
regulated version: SecurityToken deployments via \texttt{luxfi/standard}
on the Lux EVM, portfolio management via the exchange frontend
(\texttt{luxfi/exchange}), and social trading features
(leaderboards, copy-trade, rebalancing). The permission-mapping
concept was absorbed into the MPC wallet architecture: users
sign with their WebAuthn key share, the ATS co-signs after
compliance checks, and capital never leaves the user's
on-chain wallet.

\subsection{Decentralized Liquidity Pool}

\textbf{2022}: Uniswap V3--style concentrated liquidity with
up to 4000x capital efficiency, extended to support both
fungible and non-fungible assets in a single pool.

\textbf{2026}: The Lux DEX implements a native orderbook
(CLOB) rather than an AMM-only model. AMM pools exist as
fallback liquidity via \texttt{luxfi/standard/contracts/amm/}
(AMM V1--V4, Oracle-Mirrored AMM). The NFT liquidity pool
concept was not pursued --- the market moved toward
orderbook-based NFT trading (Blur/Tensor model) rather than
AMM curves for non-fungibles.

\subsection{Atomic Swaps and Bridge}

\textbf{2022}: HTLC-based atomic swaps with BTC, subsidized
by gateway deposits. Cross-chain interoperability via oracles
(Chainlink).

\textbf{2026}: The Z-Chain (zero-knowledge bridge) uses
homomorphic encryption for shielded cross-chain transfers.
The Lux DEX cross-chain gateway supports 8 chains natively.
\texttt{luxfi/standard/contracts/bridge/} implements the
Warp Messenger protocol for L1$\leftrightarrow$L2 asset
movement. Chainlink oracles replaced with Lux Oracle
(native, threshold-signed price feeds).

\subsection{Governance and DAO}

\textbf{2022}: LUX token governance with quadratic voting,
holographic consensus, proof-of-humanity biometric
authentication, proxy voting, and Ricardian contracts.

\textbf{2026}: DAO governance contracts exist in
\texttt{luxfi/standard/contracts/dao/} and
\texttt{contracts/governance/}. Voting is on-chain with
LUX staking. Quadratic voting and proof-of-humanity were
not implemented --- standard token-weighted voting with
delegation is the current model. Ricardian contracts were
superseded by on-chain compliance gates
(\texttt{IComplianceGate} interface).

\subsection{Security}

\textbf{2022}: SHA-3, PKI, time-delayed multisig, benign
fill, hardware wallets for all signing, homomorphic
zero-knowledge file systems.

\textbf{2026}: All realized plus post-quantum upgrades:
\begin{itemize}[nosep]
  \item SHA-3 $\rightarrow$ used throughout (Keccak-256 for
        EVM compatibility, SHA3-256 for non-EVM).
  \item Multisig $\rightarrow$ 2-of-3 MPC
        (CGGMP21 + FROST), strictly superior to multisig
        (single on-chain address, no $n$-of-$m$ script
        complexity).
  \item Hardware wallets $\rightarrow$ HSM co-signing
        (AWS KMS, GCP, Zymbit) for settlement intents.
  \item Zero-knowledge file system $\rightarrow$
        \texttt{PrivateStore.sol} + \texttt{/v1/base/private}
        (E2E encrypted blind blob store, shipped Q1 2026).
  \item FHE for front-running prevention $\rightarrow$ FHE
        CRDT merge (shipped Q1 2026), encrypted orderbook
        planned for Lux DEX V2.
\end{itemize}

%% ============================================================
\section{Tokenomics: 2022 Model vs.\ Current}

The 2022 paper proposed an equation-of-exchange valuation model
(Burniske: $MV = PQ$; Buterin: $MC = TH$) with an 8\% initial
inflation decreasing 15\% annually to a 1.5\% floor, and
staking rewards distributed proportionally to validators.

The current Lux Network uses a fixed supply of 10B LUX with no
inflation. Validator rewards come from transaction fees, not
token emission. The equation-of-exchange model was not adopted
--- the market moved toward fee-burn models (EIP-1559 style)
rather than inflation-funded staking.

The RWA-backed stable currencies (LUXG gold-pegged, LUXU
uranium-pegged) described in the 2022 paper remain a design
option; deployed stablecoins are USD-backed via registered
custodians. (GLUX was the Lux governance token briefly and
is unrelated to gold collateral.)

%% ============================================================
\section{Features: Shipped, Evolved, Dropped}

\begin{center}
\small
\begin{tabular}{p{4.5cm}ccp{5cm}}
\toprule
\textbf{2022 Feature} & \textbf{Status} & \textbf{Where} \\
\midrule
Non-custodial DEX         & Shipped  & \texttt{luxfi/dex} \\
Social profiles + stats   & Designed & Exchange v2 milestone \\
Copy trading              & Designed & Social trading milestone \\
People-Based Portfolios   & Designed & Portfolio rebalancing milestone \\
Trading charts            & Shipped  & Exchange frontend \\
Indicator alarms          & Partial  & Webhook events via Hanzo Tasks \\
Smart search              & Shipped  & Explorer + exchange search \\
Watchlist                 & Shipped  & \texttt{PrivateStore} (E2E encrypted) \\
Content rewards (Steemit) & Dropped  & No ROI; social features via standard UX \\
Trading bots              & Shipped  & CEX multi-protocol (FIX, WS, ZAP, JSON-RPC) \\
Token-curated support     & Dropped  & Standard support model \\
ETF baskets               & Partial  & SecurityToken framework supports basket creation \\
Atomic swaps              & Evolved  & Cross-chain gateway (8 chains) \\
Gold-backed stablecoin    & Evolved  & USDL (USD-backed, not gold) \\
RWA tokens                & Shipped  & SecurityToken framework in \texttt{luxfi/standard} \\
Gateway custodian         & Evolved  & MPC wallets replace gateway model \\
Hardware wallet (Yubikey)  & Evolved  & WebAuthn/passkeys (FIDO2 successor) \\
DAO governance            & Shipped  & \texttt{luxfi/standard/contracts/dao/} \\
Quadratic voting          & Dropped  & Token-weighted + delegation \\
Ricardian contracts       & Evolved  & \texttt{IComplianceGate} on-chain \\
IEC 61508 / CERT QA       & Partial  & CI linting + adversarial review (Red/Blue) \\
\bottomrule
\end{tabular}
\end{center}

%% ============================================================
\section{The Centralization Problem: 2022 to 2026}

The 2022 paper catalogued exchange hacks totaling \$40B+ in
stolen BTC. The thesis---that centralized custody is the root
vulnerability---remains valid. Since 2022, the collapse of FTX
(November 2022, \$8B+ in customer losses) dramatically reinforced
this point.

Lux's answer is unchanged: non-custodial MPC wallets where the
user holds one key share, the platform co-signs after compliance,
and a cold backup enables recovery. No single party can move funds
unilaterally. The PostgreSQL ledger is source of truth for trade
matching; the chain is source of truth for settlement. If the chain
goes down, trades continue; settlement catches up.

%% ============================================================
\section{Remaining Roadmap}

\begin{enumerate}[nosep]
  \item \textbf{Encrypted orderbook}: FHE-encrypted limit orders
        where the matching engine operates on ciphertexts. Eliminates
        front-running at the protocol level. Depends on FHE throughput
        improvements (LP-203, GPU NTT).
  \item \textbf{Social trading V1}: copy-trade, leaderboards,
        people-based portfolios. Architecture approved; implementation
        scoped for Q2 2026.
  \item \textbf{Cross-chain DEX aggregation}: route orders to external
        venues via the Lux Broker (16 providers) when on-chain
        liquidity is insufficient. FINRA 5310 best-execution
        compliance.
  \item \textbf{Institutional protocols}: FIX 4.4 gateway (port 8094
        on CEX), ZAP binary protocol for internal HFT, post-trade
        surveillance (wash trading, structuring, velocity detection).
\end{enumerate}

%% ============================================================
\section{Conclusion}

The 2022 LX whitepaper proposed a decentralized social trading
exchange. Four years later, the core thesis---non-custodial,
compliant, high-performance, interoperable---is realized across
four production systems: the Lux DEX (434M ops/sec native
orderbook), the Lux CEX (institutional compliance gateway), the
Lux Broker (16-venue smart order router), and the Lux ATS
(\texttt{luxfi/ats}). The social trading
features (copy-trade, PBPs, leaderboards) remain the primary
near-term milestone. The encrypted orderbook via FHE is the
primary long-term differentiator.

\bibliographystyle{plain}
\begin{thebibliography}{99}

\bibitem{uniswapv3}
H.~Adams, N.~Zinsmeister, M.~Salem, R.~Keefer, and D.~Robinson,
``Uniswap v3 Core,'' 2021.

\bibitem{nielsen2006}
J.~Nielsen, ``F-Shaped Pattern For Reading Web Content,''
Nielsen Norman Group, 2006.

\bibitem{cggmp21}
R.~Canetti, R.~Gennaro, S.~Goldfeder, N.~Makriyannis, and
U.~Peled, ``UC Non-Interactive, Proactive, Threshold ECDSA
with Identifiable Aborts,'' \textit{ACM CCS}, 2020.

\bibitem{ftx-collapse}
J.~Faux, ``The Crypto Geniuses Who Vaporized a Trillion
Dollars,'' Bloomberg, 2022.

\bibitem{atomic-swaps}
M.~Herlihy, ``Atomic Cross-Chain Swaps,'' \textit{PODC}, 2018.

\end{thebibliography}

\end{document}
